% !TEX root = ../dissertationOG.tex

\clearpage % clear the prior chapter's page
\chapter{Render-Guided Volumetric Morphing with Differentiable Physics}\label{CH4_PhysMorphGS}
% Packages used by this chapter: amsmath, graphicx, booktabs, tikz, natbib (\citep).
% The TikZ libraries below are standard; loading them here, outside any group, is safe in a chapter file,
% but the dissertation preamble is the better place for this line.
\usetikzlibrary{positioning,arrows.meta,fit,calc}

% ============================================================================
\section{Overview}
\label{sec:overview}
% ============================================================================

Chapter~\ref{CH3_SHAPEMORPHING} morphed one shape into another by controlling a differentiable material point
method (MPM) simulation through per-particle increments of the deformation gradient. The control enters the
simulation only as a stress, so the morph conserves mass and momentum and can change topology, but it was scored
by a mass residual on a grid, and its appearance was obtained afterwards by surface extraction, outside the
optimisation. This chapter closes the remaining gap
between the physics and its picture. We optimise the same kind of physical body so that it also \emph{looks}
like the target: differentiable rendering terms, evaluated on the body's own surface, steer the simulation through
its controls, while the physics keeps its conservation properties.

An earlier version of this idea, our PhysMorph-GS preprint~\citep{song2025physmorphgs}, coupled the MPM body to
Gaussian splats through a surface-aware upsampling bridge and adaptive covariances. Building on it, we found that
the coupling, the objective and the control all had to change before the render term could act on what a viewer
sees without breaking the physics (Section~\ref{ch4:sec_from_earlier}). The method of this chapter is the result.
It keeps the stress control of Chapter~\ref{CH3_SHAPEMORPHING} and the render-guided idea of the preprint, and
replaces almost everything around them:
\begin{itemize}
  \item \textbf{Settled windows.} The body is optimised window by window. Each window runs $T=20$ driven steps
        with the controls and $T=20$ released steps without them, and every term is evaluated on the released
        end state. A control is therefore scored by where the body comes to rest once it lets go
        (Section~\ref{ch4:sec_windows}).
  \item \textbf{Two controls on two scales.} The capped stress increment $\delta\mathbf{F}_c$ acts through the grid,
        as in Chapter~\ref{CH3_SHAPEMORPHING}. A normal offset $u$ of the outer particle layer, free of the six
        rigid modes, acts below the grid, where the rendering asks for changes the grid cannot represent
        (Section~\ref{ch4:sec_controls}).
  \item \textbf{A volume-conserving body.} Each particle carries the volume its motion makes, the source and
        target samples are matched in density, and the plastic assimilation takes the volume the end keeps
        (Section~\ref{sec:mpm}).
  \item \textbf{An eight-term objective.} A debiased Sinkhorn transport energy with a surface-proximity term and the
        released end's drift, the released motion, two local cleanup terms, and the silhouette and shading of the
        body's exterior surface seen from 18 views (Section~\ref{ch4:sec_objective}).
  \item \textbf{A measured render influence.} The render weight is recalibrated in every window so that the render
        term holds one third of the first gradient. Its effect is always reported against a $\lambda=0$ twin:
        the same code, seed and particles with the render weight set to zero (Section~\ref{sec:results}).
\end{itemize}
Figure~\ref{fig:overview} summarises the pipeline. The same body also takes part in physics that a geometric
interpolant cannot express: two thrown spheres collide, part in the air, and each becomes a bunny that lands and
tumbles (Section~\ref{ch4:sec_scene}).

\begin{figure*}[t]
    \centering
    \resizebox{\textwidth}{!}{\input{figs/ch4/pipeline_overview.tikz}}
    \caption{The pipeline of this chapter. Top: a morph (sphere to bunny, 300k particles) as a sequence of
    committed windows. Each window simulates $T$ driven and $T$ released steps of a differentiable MLS-MPM body (1).
    The physics loss is evaluated on the released end state (transport, surface proximity, end drift and released
    motion). The exterior surface of the particles (2) is drawn from 18 views, and its silhouette and shading are
    compared with the target's pictures (3). The physics and render gradients are taken through the same adjoint;
    the render gradient loses any component that opposes the physics gradient and is weighted so that it holds one
    third of the first gradient. After the window's iterations the state is committed, its elastic stretch is
    partly made plastic, and the next window starts from it.}
    \label{fig:overview}
\end{figure*}

% ============================================================================
\section{From the Earlier Formulation to the Present One}
\label{ch4:sec_from_earlier}
% ============================================================================

The earlier formulation~\citep{song2025physmorphgs} optimised per-particle control deformation gradients with Adam
and Lion, scored the body with the log-grid mass loss of Chapter~\ref{CH3_SHAPEMORPHING}~\citep{11088224}, and
rendered it through a six-stage upsampling bridge (surface detection, density weighting, Gumbel--Softmax importance
sampling, Taubin smoothing, adaptive covariances) into dense, render-only Gaussians whose covariances followed a
smoothed deformation gradient. Rendering gradients from several passes were cached and injected into later passes
as a stale but stable guide. Its released implementation optimised the body window by window: each window
simulated $T$ driven steps and scored the body at their end against a paced intermediate target, which moved every
particle one loss cell further along a transport plan toward the mesh. Pursuing that design to high particle counts
exposed four problems, each measured before it was changed.

\paragraph{Scoring while pushing.} The windowed method scored the body at the end of its driven steps, while the
control was still pushing. A window could accept a state that sprang back once the next window began, so the outer
layer reversed direction from one window to the next: at 300k particles, 70\,\% of the windows reversed, and every
window of the second half did. Scoring the released end instead charges the rebound to the control that caused it.

\paragraph{A loss blind to translation.} A cell-wise mass residual has its gradient only where the density jumps at
the surface; it cannot see a filled region that is merely in the wrong place. Driven toward a paced intermediate
target, it stripped the outermost particle layers ahead of the bulk, and thin features grew as a sparse haze. A
transport divergence to the fixed target replaces it: its gradient is a transport field over the whole body.

\paragraph{Rendering below the grid.} The render covector lives on the outer layer of the body (85--99\,\% of its
squared norm in our measurements), and much of it varies on the scale of a few particle spacings. After the MPM adjoint, every control
gradient has the grid's correlation length: the B-spline transfer averages a per-particle request over the
many particles that share a cell, and a sub-cell bump is not a momentum mode of the discretisation. No
upsampling of render-only samples changes this, because those samples cannot move the body. The present method
gives the render term a position channel below the grid, the layer offset $u$, and reads the render terms on a
surface built from the particles themselves.

\paragraph{Fixes outside the physics.} An intermediate extension of the windowed method stopped the reversals with
per-particle pins and commit-time position corrections. Those edits moved the body outside the simulation: its
centre of mass drifted by 0.6--0.8 particle spacings over a morph. The present method has no pins and no
commit-time position edits. The only position updates below the grid are the layer relaxation, $u$, a minimum
spacing and the re-coupling of detached fragments, the first three free of rigid motion by construction
(Section~\ref{ch4:sec_conservation}).

The settled-window design that resulted reached, on the same 300k bunny and seed, a silhouette IoU of 0.9851
against 0.9769 for the extended windowed method and 0.9646 for the earlier formulation's code, in 7 minutes
against 12 and 28. The rest of this chapter describes the formulation in its final, frozen form, and the evidence
for it.

% ============================================================================
\section{The Simulated Body}
\label{sec:mpm}
% ============================================================================

\subsection{Discretisation}
The source (a sphere) and the target mesh are each sampled into $N$ particles by stratified volume sampling, one
jittered particle per voxel of the sampler's fill. The source sample is rescaled so that its fill volume equals the
target sample's; a volume-conserving body can then fill the target exactly. Every particle has the same dynamic
mass $m_p=\bar m/N$, so the body's mass does not depend on $N$. We call $a=(V/N)^{1/3}$, the edge of the volume
per particle, the \emph{sampling pitch}; all small lengths below are given in pitches.

We simulate the body with the moving least squares material point method~\citep{hu2018mlsmpm} (MLS-MPM) with
cubic B-spline weights and APIC transfers~\citep{jiang2016material}. The MPM cell is set by the shape, not by $N$:
$\Delta x$ is the source's bounding-box diagonal divided by 26, which puts about 23 particles in a cell at
$N=40{,}000$ and about 183 at $N=300{,}000$. The time step is $\Delta t=1/240$\,s. The elastic response is fixed
corotated~\citep{2012-FixedCoratedElasticty,stomakhin2013material},
\begin{equation}
\label{eq:fc_stress}
\mathbf{P}(\mathbf{F})=2\mu\,(\mathbf{F}-\mathbf{R})+\lambda_{L}\,(J-1)\,J\,\mathbf{F}^{-\top},
\qquad J=\det\mathbf{F},
\end{equation}
with $\mathbf{R}$ the rotation of the polar decomposition, Lam\'e parameters $\mu,\lambda_L$ from $E=1.4\times10^{5}$
and $\nu=0.2$, and $J$ floored at $10^{-6}$. We write $\lambda_L$ for the Lam\'e parameter to keep $\lambda$ for
the render weight.

\subsection{One step}
Each particle carries its position $\mathbf{x}_p$, velocity $\mathbf{v}_p$, APIC affine matrix $\mathbf{C}_p$,
deformation gradient $\mathbf{F}_p$, plastic deformation gradient $\mathbf{F}^{\mathrm{P}}_p$, carried volume ratio
$J_p$ and rest volume $V^0_p$ (from the initial density). In a driven step the control $\delta\mathbf{F}_{c,p}$
enters the elastic deformation
\begin{equation}
\label{eq:control_F}
\mathbf{F}^{\mathrm{E}}_p=\big(\mathbf{F}_p+\delta\mathbf{F}_{c,p}\big)\big(\mathbf{F}^{\mathrm{P}}_p\big)^{-1},
\qquad
\boldsymbol{\tau}_p=\mathbf{P}\big(\mathbf{F}^{\mathrm{E}}_p\big)\,\mathbf{F}^{\mathrm{E}\top}_p ,
\end{equation}
and the Kirchhoff stress $\boldsymbol{\tau}_p$ enters the particle-to-grid transfer,
\begin{equation}
\label{ch4:eq_p2g}
m_i\mathbf{v}_i=\sum_p w_{ip}\Big[m_p(1-\kappa\Delta t)\,\mathbf{v}_p
  +\Big(m_p\mathbf{C}_p-\tfrac{3}{\Delta x^{2}}\Delta t\,V^0_p\,\boldsymbol{\tau}_p\Big)(\mathbf{x}_i-\mathbf{x}_p)\Big],
\qquad m_i=\sum_p w_{ip}m_p ,
\end{equation}
where $w_{ip}$ is the cubic B-spline weight and $\kappa=0.9\,\mathrm{s}^{-1}$ a weak ambient drag. After the grid
update (external forces and the walls of the domain, when present), the grid-to-particle transfer gathers
$\mathbf{v}_p=\sum_i w_{ip}\mathbf{v}_i$ and
$\mathbf{C}_p=\tfrac{3}{\Delta x^2}\sum_i w_{ip}\mathbf{v}_i(\mathbf{x}_i-\mathbf{x}_p)^{\top}$.

\paragraph{Deformation and volume.} The deformation gradient is advected with the control included and then
relaxed toward its previous value, as in the implementation of Chapter~\ref{CH3_SHAPEMORPHING}:
\begin{equation}
\hat{\mathbf{F}}_p=(\mathbf{I}+\Delta t\,\mathbf{C}_p)\big(\mathbf{F}_p+\delta\mathbf{F}_{c,p}\big),\qquad
\tilde{\mathbf{F}}_p=(1-s)\,\hat{\mathbf{F}}_p+s\,\mathbf{F}_p ,\qquad s=0.955 .
\end{equation}
This smoothing keeps 4.5\,\% of each step's increment. Applied to the volume as well, it let thin parts in transit
stream out as a dilated sheet whose stress read almost no volume change. We therefore carry the motion's own volume
ratio separately and rescale the smoothed gradient to it after every step,
\begin{equation}
\label{ch4:eq_volume}
J_p\leftarrow J_p\det(\mathbf{I}+\Delta t\,\mathbf{C}_p),\qquad
\mathbf{F}_p\leftarrow\big(J_p/\det\tilde{\mathbf{F}}_p\big)^{1/3}\,\tilde{\mathbf{F}}_p ,
\end{equation}
so that $\det\mathbf{F}_p=J_p$ while the isochoric part keeps the smoothing's shape.

\paragraph{Plastic assimilation.} When a window is committed, half of each particle's elastic stretch becomes
plastic: with $\mathbf{F}^{\mathrm{E}}=\mathbf{R}_e\mathbf{S}_e$,
$\mathbf{F}^{\mathrm{P}}\leftarrow\mathbf{S}_e^{\eta}\mathbf{F}^{\mathrm{P}}$ with $\eta=0.5$, and every principal
stretch of $\mathbf{F}^{\mathrm{P}}$ is kept in $[0.2,5]$. The assimilation takes the volumetric part of the stretch
as well, so a volume the end state keeps becomes the body's rest volume instead of being held by the control. The
end of a morph is then the body's own rest: releasing the controls changes almost nothing.

\subsection{Position updates below the grid}
\label{ch4:sec_below_grid}
The grid sees a cell, not the arrangement of the particles inside it. Three position updates act at the particle
scale, each spread over the steps of a window like a constraint projection.
\begin{itemize}
  \item \textbf{Minimum spacing.} Each particle is moved away from its neighbours that are nearer than
        $r_{\min}=0.9\,a$ (the spacing of a Poisson-disk sample of the same density), by half the overlap over one
        window:
        $\mathbf{x}_p\mathrel{+}=\tfrac{1}{2T}\sum_{q\in\mathcal{N}_p}\max(0,r_{\min}-\ell_{pq})\,
        (\mathbf{x}_p-\mathbf{x}_q)/\ell_{pq}$, $\ell_{pq}=\lVert\mathbf{x}_p-\mathbf{x}_q\rVert$.
        The neighbour lists are the symmetric closure of each particle's 16 nearest neighbours at the window's
        start, so every push is pairwise, equal and opposite, and central: it moves neither the centre of mass nor
        any rigid mode.
  \item \textbf{Outer-layer relaxation.} The outer layer (particles whose offset from the centroid of their 32
        nearest neighbours exceeds half a spacing) is relaxed toward the plane of its same-side neighbours. With
        $d_p=\mathbf{n}_p\cdot(\mathbf{x}_p-\sum_q w_{pq}\mathbf{x}_q)$ the plane residual over 24 same-side neighbours
        and $\bar d_p=\sum_q w_{pq}d_q$, each driven step moves
        $s_p=-\tfrac{1}{T}\,(d_p-\bar d_p-r_p)$ along $\mathbf{n}_p$. The reference $r_p$ is the same operator applied
        to the feet of the layer particles on the target mesh (for particles within one spacing of it), so the layer
        keeps the target's relief below a cell and loses only the sampling noise. The step's components along the six rigid modes of the layer are removed.
  \item \textbf{Fragments.} A particle that no longer shares grid support with the body (no other particle in the
        $3^3$ cells around it) takes the mean velocity of its eight source neighbours and is pulled back toward their
        rest distances over one window.
\end{itemize}
The relaxation acts only in the driven half of a window, the steps in which $u$ acts; in the released half the outer
layer is left to the physics, and the only position updates there are the minimum spacing and the fragments' bonds.

% ============================================================================
\section{Controls}
\label{ch4:sec_controls}
% ============================================================================

\paragraph{The stress increment $\delta\mathbf{F}_c$.} In each of the $T$ driven steps every particle has a
$3\times3$ increment $\delta\mathbf{F}^{n}_{c,p}$, capped per particle and step at
$\lVert\delta\mathbf{F}^{n}_{c,p}\rVert_F\le0.02$. It enters only through Eq.~\eqref{eq:control_F}. Because the stress
enters the grid as $\boldsymbol{\tau}_p(\mathbf{x}_i-\mathbf{x}_p)$ and the B-spline weights satisfy
$\sum_i w_{ip}(\mathbf{x}_i-\mathbf{x}_p)=\mathbf{0}$, the control cannot change the total linear momentum, and since
$\boldsymbol{\tau}_p$ is symmetric, APIC keeps the angular momentum as well~\citep{jiang2016material}. A window
starts from half of the previous window's control.

\paragraph{The layer offset $u$.} Every outer-layer particle has one scalar $u_p$ per window, a displacement along its
normal applied in $T$ equal parts over the driven steps:
\begin{equation}
\label{ch4:eq_u}
\mathbf{x}_p\leftarrow\mathbf{x}_p+\Big(\tfrac{1}{T}\,g_p\,u_p+s'_p\Big)\mathbf{n}_p ,
\end{equation}
where $s'_p$ is the rigid-free relaxation step above and $g_p\in\{0,1\}$ a gate. Its adjoint is the identity times
the physics' response, so the render covector reaches it without the grid's low-pass. Three rules keep it a surface
control and not a hand edit:
\begin{itemize}
  \item it is gated: $u$ acts only on layer particles whose remaining transport, the label-free transport map of
        the window's start on the MPM grid, is within one MPM cell, so it shapes an arriving surface rather than
        carrying material;
  \item it is projected free of the six rigid modes of the layer after every step of the optimiser, so it moves
        the body neither along nor about any axis;
  \item it is bounded by one spacing per window.
\end{itemize}
What $u$ moves is held by the released physics or not; the released half decides.

% ============================================================================
\section{Settled Windows}
\label{ch4:sec_windows}
% ============================================================================

A window simulates $2T=40$ steps ($1/6$\,s): $T$ driven steps with the controls, then $T$ released steps with the
controls set to zero while the physics keeps running. Every term of the objective is evaluated on the released end
state $(\mathbf{x}_{2T},\mathbf{v}_{2T})$ and on the released velocities. Rebound, ringing and anything else the
control holds only while it pushes therefore cost the control directly. The optimiser runs eight iterations per
window; the window is then committed (Section~\ref{sec:optimization}) and the next window starts from the committed
state. A 300k bunny morph takes 76--106 committed windows, a dragon 139--186.

% ============================================================================
\section{Objective}
\label{ch4:sec_objective}
% ============================================================================

The objective has eight terms in three groups,
\begin{equation}
\label{ch4:eq_objective}
E=\underbrace{\sigma\big(S_\varepsilon+E_{\mathrm{prox}}+E_{\mathrm{drift}}\big)+E_{\mathrm{rel}}}_{E_{\mathrm{phys}}}
 \;+\;\underbrace{E_{\mathrm{spray}}+E_{\mathrm{near}}}_{E_{\mathrm{clean}}}
 \;+\;\lambda\underbrace{\big(E_{\mathrm{sil}}+E_{\mathrm{shade}}\big)}_{E_{\mathrm{render}}} .
\end{equation}
Which terms are needed was decided by switching terms off on a 19-mesh gallery at $N=40{,}000$ and at 300k. The
geometry and settling terms carry no weight of their own: $\sigma$ is one constant measured once at the source, which
gives the geometry energy, there, the gradient norm of a log-mass grid loss of the kind used in
Chapter~\ref{CH3_SHAPEMORPHING}, and so fixes the unit of the two cleanup weights (0.2 each). There is no regulariser of the control or of the volume: the earlier
control magnitude and smoothness terms were $10^{-10}$ and $10^{-8}$ of the objective, and removing them together
with a volume prior changed none of the measured quantities (det\,$\mathbf{F}$, anisotropy, control and surface
roughness). The cap on $\delta\mathbf{F}_c$, the released half and a validity check of the rollout bound the
controls instead.

\subsection{Physics core}
\paragraph{Transport.} The body's mass and the target sample's are rasterised on a loss grid, and $S_\varepsilon$ is
their debiased Sinkhorn divergence~\citep{feydy2019sinkhorn},
\begin{equation}
S_\varepsilon(\alpha,\beta)=\mathrm{OT}_\varepsilon(\alpha,\beta)-\tfrac12\mathrm{OT}_\varepsilon(\alpha,\alpha)
 -\tfrac12\mathrm{OT}_\varepsilon(\beta,\beta),\quad
\mathrm{OT}_\varepsilon(\alpha,\beta)=\min_{\pi\in\Pi(\alpha,\beta)}\langle\mathbf{c},\pi\rangle
 +\varepsilon\,\mathrm{KL}(\pi\,\Vert\,\alpha\otimes\beta),
\end{equation}
with the squared Euclidean cost and blur $\varepsilon=\Delta x_\ell^2$, one loss cell. The loss cell equals the MPM
cell up to $N=40{,}000$ and follows the particle spacing above it,
$\Delta x_\ell=\Delta x\,(40{,}000/N)^{1/3}$ (about the bounding-box diagonal over 51 at 300k). The problem is solved
in the log domain with $\varepsilon$-scaling from the squared grid diameter down to the blur. The cost is separable by
axis, so every Sinkhorn sweep is three one-dimensional log-sum-exp passes on the GPU. Its gradient follows from the
envelope theorem: the converged potentials, differentiated through the rasterisation weights. The transport moves
mass and is indifferent to topology; unlike a mass residual, its gradient is a transport field over the whole body.

\paragraph{Surface proximity.} The transport's blur cannot tell apart positions within a loss cell. The proximity
term asks, at every outer point $\mathbf{y}$ of the target sample, whether body material is there:
\begin{equation}
E_{\mathrm{prox}}=\rho^2\,\operatorname*{mean}_{\mathbf{y}}\;
 \max\!\Big(0,\,1-\frac{K(d_{\min}(\mathbf{y}))}{\tfrac12K(\varsigma)}\Big)^{2},
 \qquad K(d)=e^{-d^2/2h^2},
\end{equation}
where $d_{\min}(\mathbf{y})$ is the distance to the nearest body particle, $\rho$ the target sample's median distance
to its 8th neighbour, $h=\rho/2$, and $\varsigma$ its median nearest-neighbour spacing. The term is zero where a
particle sits within about $1.5$ spacings of $\mathbf{y}$ and largest at a gap several spacings wide; it has no
weight and no bound.

\paragraph{Settling.} Two terms ask for an end at rest:
\begin{equation}
E_{\mathrm{drift}}=(T\Delta t)^2\operatorname*{mean}_p\lVert\mathbf{v}_p(t_{2T})\rVert^2 ,\qquad
E_{\mathrm{rel}}=(T\Delta t)^2\operatorname*{mean}_{t\in\text{released},\,p}\lVert\mathbf{v}_p(t)\rVert^2 .
\end{equation}
The end drift is the squared displacement the end velocity would add over one more horizon and belongs to the
geometry energy. The released motion is the stability term: without it, several gallery meshes stopped after 8--11
windows while still moving; without the end drift, runs were sound but their released end moved 1.4--2.9 times
faster.

\subsection{Cleanup}
Two local terms with fixed weights act on particles the transport cannot resolve.
The \emph{spray pull} is a one-sided Wasserstein-1 pull, $\sum_p g_p m_p\,d_{\mathcal{T}}(\mathbf{x}_p)$, down the
target's distance field $d_{\mathcal{T}}$, on the particles an isolation gate $g_p$ marks (a ramp of the distance to
the 8th neighbour from 1.2 to 1.8 times its median, frozen per window). The \emph{near band} pulls each particle that
lies between the sampling berth (about two target spacings) and one loss cell from the target toward its nearest
target point, $\sum_p e_p m_p\max(0,\lVert\mathbf{x}_p-\mathbf{y}_{a(p)}\rVert-b)$, where the transport's blur cannot
tell positions apart.

\subsection{Render terms on the exterior}
\label{ch4:sec_exterior}
The render terms read the surface of the body, not the particle cloud. The body's \emph{exterior} is the zero set of
Zhu and Bridson's field~\citep{zhu2005sand},
\begin{equation}
\phi(\mathbf{q})=\lVert\mathbf{q}-\bar{\mathbf{x}}(\mathbf{q})\rVert-r_o,\qquad
\bar{\mathbf{x}}(\mathbf{q})=\frac{\sum_p k_p(\mathbf{q})\,\mathbf{x}_p}{\sum_p k_p(\mathbf{q})},\quad
k_p(\mathbf{q})=\max\!\Big(0,1-\tfrac{\lVert\mathbf{q}-\mathbf{x}_p\rVert^2}{R^2}\Big)^{3},
\end{equation}
with kernel radius $R=3a$ and offset $r_o=0.8a$. Surface discs $(\mathbf{p}_j,\mathbf{n}_j)$ sit one to each cell of a
lattice fixed in space that the zero set crosses (the cell centre projected onto the zero set), at a spacing of
$\min(\text{half a picture pixel},\,0.92a)$. The discs follow the particles between searches by one Newton step of the
field and are found again when a tenth of them has moved more than half a lattice spacing; because the lattice is
fixed, the same surface yields the same discs.

The pictures are 18 orthographic views (six azimuths at elevations $0$ and $\pm0.5$\,rad) at 96\,px for 40k
particles, a resolution that follows $N$ as $96\,(N/40{,}000)^{1/3}$ (188\,px at 300k). The silhouette splats the discs
with opacity $\alpha=1-e^{-k\,w}$ ($k=1.5$, $w$ the bilinear coverage) and penalises missing coverage inside the
target more than excess outside it,
\begin{equation}
E_{\mathrm{sil}}=\operatorname*{mean}_{\text{views, pixels}}\Big[w_h\max(0,\alpha^\ast-\alpha)^2
 +w_s\max(0,\alpha-\alpha^\ast)^2\Big],\qquad w_h=2,\;w_s=1 .
\end{equation}
The shading term compares headlight-Lambertian images, $b=0.25+0.75\max(0,\mathbf{n}_j\cdot\mathbf{l})$, each disc
weighted by its facing cosine and a soft front bias and masked by the coverage:
$E_{\mathrm{shade}}=\operatorname{mean}(s-s^\ast)^2$. The target pictures $\alpha^\ast,s^\ast$ are drawn by the same
operator from the discs of the target sample's exterior, averaged over eight independent samples of the mesh, so
the reference carries no single sample's noise.

\paragraph{Why the exterior.} Read on the particle cloud, the render terms improve the cloud's own silhouette but
leave the surface a viewer sees where the $\lambda=0$ twin has it. Read on the exterior, they act on the displayed
surface; the comparison is in Section~\ref{ch4:sec_ablation}.

\paragraph{Display.} The images in this chapter draw the same exterior with surface discs in a studio look. The
display uses the field with Solenthaler et al.'s correction~\citep{solenthaler2007unified}, which differs from Zhu
and Bridson's only in concave gaps (0.4--5.2\,\% of the discs move, by a median $10^{-4}$ pitches) and removes blobs the
uncorrected field draws in the empty space between near surfaces, such as the dragon's mouth. Every image is a
committed simulated state; every metric is computed on the particles or on the uncorrected field.

% ============================================================================
\section{Optimisation}
\label{sec:optimization}
% ============================================================================

\paragraph{Gradients.} The physics core, the cleanup and the render terms are differentiated separately through the
same $2T$-step adjoint back to every $\delta\mathbf{F}^n_{c}$ and to $u$. The simulation and its adjoint are recorded on
a tape and captured as CUDA graphs. Finite differences agree with the adjoint within 3\,\% for every term at windows 1
and 20 of a 300k run.

\paragraph{Projection and weight.} Let $\mathbf{g}_p$ be the gradient of the physics core and $\mathbf{g}_r$ that of the
render term, each over all control leaves. The render gradient loses its component against the physics gradient
when the two conflict~\citep{yu2020pcgrad},
\begin{equation}
\tilde{\mathbf{g}}_r=\mathbf{g}_r-\frac{\min(0,\langle\mathbf{g}_p,\mathbf{g}_r\rangle)}{\lVert\mathbf{g}_p\rVert^{2}}\,\mathbf{g}_p ,
\end{equation}
and the render weight is set at each window's first gradient by
\begin{equation}
\label{ch4:eq_lambda}
\lambda\,\lVert\tilde{\mathbf{g}}_r\rVert=\tfrac12\,\lVert\mathbf{g}_p\rVert ,
\end{equation}
and held for the window. The render term therefore holds exactly one third of the first gradient's norm,
$\lambda\lVert\tilde{\mathbf{g}}_r\rVert/(\lVert\mathbf{g}_p\rVert+\lambda\lVert\tilde{\mathbf{g}}_r\rVert)=1/3$, in every
window. The physics gradient falls by two to three orders of magnitude while the body arrives, and $\lambda$ falls with
it. An earlier form averaged $\lambda$ across windows; it lagged that fall, kept the render at 0.46--0.70 of the step in
windows 3--11, and there drove protruding particles out of the body. The step direction is
$\mathbf{g}=\mathbf{g}_p+\lambda\tilde{\mathbf{g}}_r+\mathbf{g}_{\mathrm{clean}}$.

\paragraph{Step.} The step is Adam's~\citep{kingma2017adammethodstochasticoptimization} with a backtracking line search:
a trial is accepted only when the full objective decreases by the Armijo amount (or by more than the rollout's
replay noise, measured at each window's start) and the trajectory stays valid (finite, no inverted
$\mathbf{F}$, inside the domain). After each trial, $\delta\mathbf{F}_c$ is capped and $u$ is bounded and projected free of
rigid modes. The Adam moments of $\delta\mathbf{F}_c$ start afresh in every window. The moments of $u$ carry from one
committed window to the next, each coordinate's bias correction counting its own age. Restarted every window, $u$'s
first step was full-sized whatever its gradient and reversed from one window to the next, half of the late per-window
wobble; the relaxation pulling the layer back during the released half was the other half
(Section~\ref{ch4:sec_below_grid}).

\paragraph{Commit, acceptance and stop.} A window's result is committed by replaying the accepted candidate on the
exterior discs it was scored with; the plastic assimilation is applied; and the committed state is scored by one
selection merit, the objective at that state. A window that raises the merit by more than 5\,\% is rejected and the
previous state kept. Because the merit is linear in $\lambda$, the references are rescored whenever $\lambda$ changes.
The run ends after three consecutive rejections or five windows without a relative improvement of 0.3\,\%, and
delivers the trajectory up to its best window.

\subsection{What the body conserves}
\label{ch4:sec_conservation}
Mass is exact by construction. The stress control cannot change the linear or angular momentum. The minimum
spacing's push is pairwise and central. The relaxation and $u$ are free of the rigid modes of the layer. The
drag damps $m_p\mathbf{v}_p$ only, not APIC's affine momentum, so it partly decays a rigid spin and can trade angular
momentum between the orbital and affine parts; we report rotation as a kinematic quantity and angular momentum through
per-window records. Over whole 300k morphs the centre of mass moves by $1.6$--$6.9\times10^{-5}$ pitches at most, the
float32 floor of the grid transfers, and the body's physical volume $\sum_p J_pV^0_p$ stays within $-0.8$\,\% and
$+3.1$\,\% of the target's over the run and ends $0.1$--$0.4$\,\% above it (bunny and dragon).

% ============================================================================
\section{Multi-Body Scenes}
\label{ch4:sec_scene}
% ============================================================================

A single world-fixed objective cannot let a body fly, spin or land: the target, the cameras and the settling terms
are fixed in the world, so a falling centre of mass raises the objective whatever the control does. For scenes with
several bodies, contact and gravity we extend the single-body window with four definitions, keeping every per-body
term unchanged.
\begin{itemize}
  \item \textbf{One grid field per body.} On one grid, every node inside both kernels carries one velocity, so two
        bodies within two cells of each other move as one and stick. Each body deposits into its own grid field. At a
        node where both fields have mass, the bodies are in contact when their node volumes fill the cell,
        $V_0+V_1\ge\Delta x^3$ (the touching condition for any partition-of-unity kernel), and approach,
        $(\mathbf{v}_0-\mathbf{v}_{cm})\cdot\mathbf{n}>0$, with $\mathbf{n}$ the normalised difference of the fields' mass
        gradients. Body~0 then loses its normal velocity relative to $\mathbf{v}_{cm}$ with Coulomb friction ($\mu=0.5$),
        and body~1 takes the opposite change weighted by the masses, following the multi-velocity rule of Bardenhagen
        et al.~\citep{bardenhagen2000granular}. The changes are equal and opposite at each node, so linear and angular
        momentum are conserved; the floor is the same rule with a body of infinite mass.
  \item \textbf{A body frame.} Each body's objective is evaluated in its shape-matching frame, $\mathbf{y}=\mathbf{R}^{\top}
        (\mathbf{x}-\mathbf{c})$ with $\mathbf{c}$ the centroid and $\mathbf{R}=\operatorname{polar}\big(\sum_p(\mathbf{x}_p-\mathbf{c})
        (\mathbf{X}_p-\bar{\mathbf{X}})^{\top}\big)$, $\mathbf{X}$ the body's source sample. The polar factor is
        differentiated exactly, so the objective is invariant to each body's rigid motion and its position gradient
        carries no net force or torque. The cameras ride with the body, and the velocity terms read the deformation
        velocity $\mathbf{v}-(\mathbf{v}_c+\boldsymbol{\omega}\times\mathbf{r})$, so flight and tumbling cost nothing.
  \item \textbf{Time always advances.} Every window commits its line-searched control, or the uncontrolled continuation
        when no step was accepted; there is no rejection or stop. The ambient drag is replaced by a damping of the
        deformation velocity alone, which keeps each body's momentum.
  \item \textbf{A moving target.} Each body's target at time $t$ is the displacement interpolation of its transport
        read from the target's side, $\mathbf{Z}_s=\mathbf{Y}+(1-s)\,\mathbf{D}_Y(\mathbf{Y})$ with $s=3\tau^2-2\tau^3$,
        $\tau=t/t_m$, so the bodies meet half-formed instead of as finished shapes~\citep{mccann1997convexity}.
\end{itemize}
The single-body path is unchanged by these extensions, bit for bit.

% ============================================================================
\section{Implementation Details}
\label{sec:implementation}
% ============================================================================

The simulation and its adjoint are written in NVIDIA Warp. The run stage is GPU-only: particle states stay on the
device between windows, and neighbour queries and grid morphology run on the GPU through a device KD-tree; only the
preparation (mesh loading and sampling) runs on the CPU and is cached. The P2G transfer visits particles in cell order
so that the 32 lanes of a warp share a stencil and add each node once, which removed most of the same-address atomic
contention at 183 particles per cell. Two further exact speed-ups skip the layer kernels in steps where neither the
relaxation nor $u$ acts and reuse the proximity term's neighbour lists within an iteration; they leave every result
bit for bit identical. The time of an optimiser attempt is dominated by the Sinkhorn solves, the surface proximity and
the three adjoint sweeps (Section~\ref{ch4:sec_runtime}). The deliverable videos are rendered at $3840\times2160$,
20 frames per second, from every 12th simulated step.

% ============================================================================
\section{Results}
\label{sec:results}
% ============================================================================

\subsection{Protocol}
\label{ch4:sec_protocol}
Every morph starts from a sphere and targets one of six meshes at $N=300{,}000$ (bunny, dragon, teapot, bob, armadillo,
homer); a 19-mesh gallery runs at $N=40{,}000$. Three of the six main targets are not of genus zero (the dragon has
genus~15; the teapot and bob have genus~1). Runs use seeds 97, 99 and 101, and the bunny and dragon also 103 and 105.

\paragraph{Independent sample.} No result is scored against the sample it was optimised toward. Every end metric
compares the result with an \emph{independent} stratified sample of the same mesh, drawn with a seed no run uses for
its target or its pictures. The \emph{sampling floor} is the same metric between the run's own target sample and the
independent one: what a perfect placement of the particles would score.

\paragraph{Measures.}
\begin{itemize}
  \item \emph{Particle silhouette IoU} over 24 views at 128\,px (a binary $3\times3$-pixel point splat) and the
        symmetric \emph{chamfer} distance.
  \item \emph{Exterior silhouette and shading residuals}: the render terms' operators applied to the exterior of the end
        state and of the independent sample, 18 views at 96\,px.
  \item \emph{Displayed-surface} $1-\mathrm{IoU}$ of the base display at 4K against the independent sample's, front and
        back.
  \item \emph{Thin coverage}: the share of the target's thin outer points (local thickness below two MPM cells) with
        no particle within 1.5 spacings (against the run's own sample).
  \item \emph{At rest}: the mean squared particle speed at the last committed window's released end.
  \item \emph{Momentum}: the largest displacement of the centre of mass and the net rotation over all kept frames.
  \item \emph{Floating pieces}: sets of particles not linked to the body within 1.5 pitches; a piece is
        \emph{serious} if it has at least ten particles, lies at least five pitches from the body and persists.
\end{itemize}

\paragraph{Reading.} Every render result is paired with its $\lambda=0$ twin. We call two arms different on a measure
only when their seed ranges do not overlap, and changes are median against median (for an IoU, the change of
$1-\mathrm{IoU}$). Wall-clock times are on one RTX~6000 Ada with nothing else on the card.

\begin{figure*}[t]
    \centering
    \includegraphics[width=\textwidth]{figs/ch4/fig_morph_sequence.jpg}
    \caption{Sphere to bunny (top) and sphere to dragon (bottom), 300k particles, side view, the method as frozen
    (seed 97). Columns: $t=0$, $0.5$, $1$, $2$ and $4$\,s and the delivered end (12.85\,s and 30.2\,s). The body rises as
    one piece and the features grow out of it in the first second. Narrow gaps of the dragon, below two MPM cells, open
    late: at 4\,s the gap between hind leg and body is still opening, with beads of material across it
    (Section~\ref{ch4:sec_limitations}).}
    \label{ch4:fig_morph}
\end{figure*}

\subsection{Morphs}
Figure~\ref{ch4:fig_morph} shows the bunny and the dragon. The body rises as one piece and the features grow out of it;
both the physics and the render terms do 90\,\% of their (linear) drop in the first 0.6--1.6\,s while the body travels, the
render term 0.06--0.40\,s behind the physics because it reads the arriving surface (Figure~\ref{ch4:fig_progress}).
After 3\,s comes a long tail, 79--91\,\% of the delivered frames, carrying only 2--13\,\% of the particle movement: the
visible surface creeps by 2--4 pixels at 4K on four meshes, 4--6 on homer and 9--13 on the dragon, while both terms keep
falling on a logarithmic scale.

At the end, against an independent sample, the particle silhouette IoU is 0.9864--0.9899 on the six meshes (medians
over seeds) against sampling floors of 0.9883--0.9919, and the chamfer distance lies below the floor on every mesh
because the ends are more even than any random sample. On the run's own samples, 2.4--4.1\,\% of the target's thin outer
points are left uncovered. The ends are at rest (mean squared speed $4\times10^{-8}$ to $1.3\times10^{-5}$ in squared world units per
second squared). No serious
piece occurs in any of the 44 runs: at the end every piece off the body lies 1.7--3.4 pitches from it, the largest has
13 particles, and pieces farther than five pitches appear only between 1.6 and 5.8\,s, of two to five particles
(one 13-particle piece in one frame of a twin).

\begin{figure}[t]
    \centering
    \includegraphics[width=\linewidth]{figs/ch4/fig_progress.png}
    \caption{How the objective falls over a 300k morph (full method; seed 97 thick, other seeds thin). Left: the share of
    each term's drop already done; right: the share of its decades already fallen; dashed: the share of the total
    particle movement already done. Both the physics core and the render term do most of their drop in the first one to
    two seconds; on the logarithmic scale both keep falling to the end.}
    \label{ch4:fig_progress}
\end{figure}

\subsection{The effect of the render terms}
\label{ch4:sec_render_effect}
Table~\ref{ch4:tab_b1} compares the full method with its $\lambda=0$ twin on the six meshes. On every mesh the full
method lowers the exterior silhouette residual by 31--57\,\%, the exterior shading residual by 19--35\,\% and the
displayed surface's $1-\mathrm{IoU}$ by 18--33\,\% (front) and 16--33\,\% (back), with the seed ranges apart. Read on
every kept frame of the bunny and the dragon, the full method is never behind its twin on these measures: on the bunny
it leads from 0.1\,s on, the error ratio falling from 0.88 at 0.5\,s to 0.56 at the end; on the dragon the twin catches up
while the body spreads (ratio 0.99 at 2\,s) and the full method leads without a break from 2.8\,s, ending at 0.44. The
full method's own silhouette term ends at 0.10--0.26 times the twin's, while the physics objective does not pay for it:
transport, released motion and end drift end inside the twin's range on five meshes and below it on the teapot.

What the render terms do not move is equally clear. The particle silhouette IoU, the chamfer distance, the thin
coverage, the centre of mass and the rest at the end lie within the twin's seed range on every mesh. Two measures are
worse. Bob's net rotation rises from $0.0009^\circ$ to $0.0044^\circ$, and the angular momentum the drag creates grows
with the extra motion on bunny, bob, armadillo and homer (the grid's summed change of angular momentum $+7$ to
$+37$\,\%; the largest $\lvert\mathbf{L}\rvert$ at a window start up to $+68$\,\% on bob). The render term deforms these bodies more, and the frozen drag, which does not respect angular momentum,
acts more on a body that moves more. Homer's runs are 34\,\% longer with the render term.

\begin{table}[t]
\centering
\footnotesize
\setlength{\tabcolsep}{3.5pt}
\caption{The full method against its $\lambda=0$ twin at the end of 300k morphs (medians over seeds: bunny and dragon
five, the others three; full method / twin and the change). Exterior residuals and displayed surface against an
independent sample; $^\ast$: the seed ranges do not overlap.}
\label{ch4:tab_b1}
\resizebox{\linewidth}{!}{%
\begin{tabular}{lcccccc}
\toprule
mesh & silhouette IoU & ext.\ silhouette & ext.\ shading & display front & thin & at rest \\
 & (indep.) & ($10^{-4}$) & ($10^{-4}$) & $1-\mathrm{IoU}$ (\%) & (\%) & ($10^{-6}$) \\
\midrule
bunny    & 0.9873 / 0.9863 & 7.45 / 13.0 ($-43\%^\ast$) & 4.74 / 6.46 ($-27\%^\ast$) & 0.93 / 1.19 ($-22\%^\ast$) & 3.97 / 3.85 & 1.31 / 1.51 \\
dragon   & 0.9874 / 0.9874 & 5.39 / 12.6 ($-57\%^\ast$) & 4.64 / 7.01 ($-34\%^\ast$) & 0.92 / 1.27 ($-28\%^\ast$) & 3.26 / 3.23 & 7.67 / 7.64 \\
teapot   & 0.9899 / 0.9892 & 7.02 / 11.8 ($-40\%^\ast$) & 4.15 / 6.40 ($-35\%^\ast$) & 0.85 / 1.11 ($-23\%^\ast$) & 2.42 / 3.28 & 0.35 / 0.82 \\
bob      & 0.9895 / 0.9886 & 4.39 / 6.67 ($-34\%^\ast$) & 3.20 / 3.94 ($-19\%^\ast$) & 0.61 / 0.75 ($-18\%^\ast$) & 3.83 / 4.29 & 0.04 / 0.10$^\ast$ \\
armadillo& 0.9864 / 0.9861 & 6.21 / 12.2 ($-49\%^\ast$) & 3.76 / 5.52 ($-32\%^\ast$) & 1.21 / 1.81 ($-33\%^\ast$) & 4.09 / 3.79 & 13.2 / 10.3 \\
homer    & 0.9882 / 0.9880 & 3.33 / 4.84 ($-31\%^\ast$) & 1.99 / 2.60 ($-23\%^\ast$) & 0.76 / 1.03 ($-26\%^\ast$) & 4.04 / 4.90 & 1.95 / 2.82 \\
\bottomrule
\end{tabular}}
\end{table}

\paragraph{Why the particle IoU barely registers it.} The binary silhouette IoU splats each particle as a
$3\times3$-pixel point at 128\,px, about 1.8 pitches a pixel, and 72--89\,\% of the twin's error is the sampling floor
itself. Moving the full method's outer layer uniformly by 0.1 pitch along its normal changes $1-\mathrm{IoU}$ by
0.0008--0.0025, as much as the render term's whole effect (0.0005--0.0021); see Figure~\ref{ch4:fig_iou}. The render term
changes the exterior by fractions of a pitch, concentrated on silhouette rims and thin parts, below what the binary
splat resolves. The full method's outer layer sits 8--15\,\% nearer the independent sample than the twin's in its normal
offset, and its own silhouette term falls below what an independent sample of the mesh scores after a tenth to a half
of the run.

\begin{figure}[t]
    \centering
    \includegraphics[width=0.62\linewidth]{figs/ch4/fig_iou_sensitivity.png}
    \caption{The particle silhouette IoU (binary splat, 128\,px) against a uniform shift of the full method's outer layer
    along its normal, bunny and dragon at three seeds each. Dotted: the whole difference between the full method and its
    twin. A shift of a tenth of a pitch moves this measure as much as the render term's entire effect.}
    \label{ch4:fig_iou}
\end{figure}

\paragraph{Against the exact mesh.} An independent sample still carries sampling noise. Read instead against the exact
target mesh by ray casting, with no sample in the reference (Figure~\ref{ch4:fig_render_end}), the full method's exterior
outline deviates 14--33\,\% less than the twin's on all six meshes (outline shape error, the spread of the outline's
offset from the mesh's), and less than a perfect particle sample of the mesh, which the twin's does not. Its exact binary
silhouette error at 1024\,px is 4--12\,\% lower, its normals 3--8\,\% and its headlight shading 8--13\,\% closer. Every
particle exterior, a perfect sample included, sits 0.7--0.9 pitch outside the mesh because of the field's offset; this
shared dilation dominates raw silhouette scores against the mesh, which is why the offset-free measures carry the shape.
% TODO(results pending: standard metrics of the end states against the exact mesh, full method against its twin on the
% six 300k meshes and the 40k gallery: boundary IoU and boundary F-score, contour Chamfer, 3D F-score, Chamfer-L1,
% normal consistency, and novel-view PSNR / SSIM / LPIPS; report them here beside the measures above, or state plainly
% that none of them sees the render term.)

\begin{figure*}[t]
    \centering
    \includegraphics[width=\textwidth]{figs/ch4/fig_render_effect_end.png}
    \caption{The end of the run, full method against its $\lambda=0$ twin on six meshes (dot: median over seeds, bar:
    seed range; dashed: the twin; grey tick: a perfect particle sample of the mesh). Top five rows: against an
    independent sample of the target; bottom four: against the exact mesh. The render terms lower every measure that
    reads the surface; the particle silhouette IoU and the chamfer distance barely see the change.}
    \label{ch4:fig_render_end}
\end{figure*}

\paragraph{A causal test.} If the late render term maintains the surface it shapes, removing it should lose the lead and
adding it should build one. We switched the render term off at 6\,s in the full method and on at 6\,s in the twin
(bunny and dragon, two seeds each; Figure~\ref{ch4:fig_switch}). Switched off, the run loses half of the lead within
4--6 windows (0.7--1\,s) on the exterior silhouette and within 2--5 on the shading, and 63--92\,\% of it within ten
windows; 10--26\,\% of the silhouette lead remains, the part made by early pushes. Switched on, the twin builds half the
lead within 1--2 windows and 82--121\,\% of it after three, and ends level with the run that had the render term
throughout (1.00--1.07 times its exterior silhouette residual). The pre-registered criterion, half the lead lost or built
within ten windows and still at twenty on both primary measures for at least three of four mesh-seed pairs, holds on
all four for both arms. The switch moves no momentum: the centre of mass moves $2.6\times10^{-5}$--$1.5\times10^{-4}$
pitches and the net rotation is 0.003--0.007$^\circ$, within the range of the unswitched runs.

\begin{figure*}[t]
    \centering
    \includegraphics[width=\textwidth]{figs/ch4/fig_render_switch.png}
    \caption{The render term switched at 6\,s. Top: off from 6\,s in the full method; bottom: on from 6\,s in the twin.
    Exterior silhouette residual against an independent sample, relative to the twin's at the same time (twin = 1).
    Without the render term the lead goes within about a second; with it, the lead comes within a third of a second.}
    \label{ch4:fig_switch}
\end{figure*}

\paragraph{Where the render gradient acts.} A per-window counterfactual, the same window optimised again from the same
start with $\lambda=0$, separates what the render term changes (Figure~\ref{ch4:fig_rgv}). Late in a run, the render's
displacement of the end state is 75--77\,\% on the outer layer (5.7--8.9\,\% of the particles), 15--18 times the average
per particle on silhouette rims and 5--8 times on thin parts. On the layer, the render's share of the per-particle pull
is a median 20\,\% at window~1 on the bunny (35\,\% on the dragon) and 7--19\,\% later, with 37--44\,\% at the
90th percentile early in the run, on the ears, the face and the feet. Its directions mostly agree with the physics early (63\,\% of the layer where it
acts against 16\,\% in conflict at window~1) and conflict more late (41\,\% against 32\,\%), where the projection removes
the opposing part in 29 of 79 windows on the bunny. Most of the render's demand lies below the MPM cell: only 5--9\,\% of
it is at the cell scale, and late in the run $\delta\mathbf{F}_c$ aligns with that part only (cosine 0.58--0.66) and with
the render's wanted move as a whole at 0.12--0.14, while $u$ reaches it at 0.79--0.81. The cap on $\delta\mathbf{F}_c$ binds in none of
the late windows. Late in the run a window without the render term gives back 67--81\,\% of the previous window's
render-driven displacement: the late render term maintains a sub-pitch offset against the physics' own refit rather than
accumulating change.

\begin{figure*}[p]
    \centering
    \includegraphics[width=\textwidth]{figs/ch4/fig_render_gradient_early.jpg}\\[4pt]
    \includegraphics[width=\textwidth]{figs/ch4/fig_render_gradient_late.jpg}
    \caption{The render gradient on the deforming 300k bunny, early ($t=1.2$\,s, top) and late ($t=5.0$\,s, bottom). In
    each block: where the render acts on $\delta\mathbf{F}_c$ (its share of the pull on each particle's control, grey to
    orange); how strong its pull is against its typical value at that moment; where the render and the physics pull
    the same way (blue) or against each other (red), averaged over five windows; and what the render changed in
    $\delta\mathbf{F}_c$, from a $\lambda=0$ counterfactual of every window (green; arrows: the principal axis of the change,
    stretch or compression). Hatched: the render barely acts there.}
    \label{ch4:fig_rgv}
\end{figure*}

\paragraph{The 40k gallery.} On 19 meshes at $N=40{,}000$ and three seeds, the full method's particle silhouette IoU is
above its twin's at the three-seed mean on 18 of 19 meshes (bob is level), with all three seeds agreeing on 10; a paired
Wilcoxon test over the meshes gives $p=1.1\times10^{-5}$. The gains are small in absolute terms (median $+0.0007$), a
median quarter of the twin's gap to the sampling floor, and the lead has halved relative to the recipe before its three late
refinements (Section~\ref{ch4:sec_ablation}), because those refinements lifted the twin more than the full method.

Read against the exact meshes (Figure~\ref{ch4:fig_g40}; all 19 meshes, three seeds each), the render term at 40k
reshapes the outline rather than the silhouette's area. Its outline shape error is 3--15\,\% below the twin's on all 19
meshes (56 of 57 mesh-seed pairs, Wilcoxon $p=5\times10^{-11}$; median $-10$\,\%, about half the effect at 300k), and its
normals come about 2\,\% closer (44 of 57, $p=6\times10^{-6}$). Its exact binary silhouette error, in contrast, is about
1\,\% higher (42 of 57), because its exterior sits slightly farther outside the mesh (45 of 57); at 300k the full method's
exterior was the nearer one and its exact silhouette error 4--12\,\% lower. An area measure such as the gallery's binary
IoU therefore shows almost no lead at 40k, while the outline improves on every mesh.
% TODO(results pending: why the full method's 40k exterior sits slightly farther outside the mesh than its twin's, the
% opposite of 300k; write the finding here once read.)

\begin{figure}[t]
    \centering
    \includegraphics[width=0.85\linewidth]{figs/ch4/fig_gallery40k_mesh.png}
    \caption{The 40k gallery against the exact target meshes: the full method's error relative to its $\lambda=0$ twin
    (``arm'' in the panel titles; dot: median over three seeds, bar: range, dashed: the twin, grey tick: a perfect
    particle sample of the mesh). Left: exact binary silhouette $1-\mathrm{IoU}$ at 1024\,px; middle: outline shape error;
    right: normal angle at the nearest mesh point. The outline improves on all 19 meshes; the silhouette's area does not.}
    \label{ch4:fig_g40}
\end{figure}

\paragraph{Materials and external forces.} Every material reaches its target at 300k (particle IoU 0.9863--0.9876):
a soft solid ($E\times0.1$) arrives more slowly (transport energy at 2\,s 55--178 times the default's), a stiff one
($E\times7$, the stiffest the fixed step carries) arrives faster but ends less at rest (2.2--3.5 times the default's
end kinetic measure), and $\nu=0.45$ costs 9--21\,\% on some surface and thin measures. The render term's lead holds in every material: exterior silhouette
$-29$ to $-57$\,\%, shading $-20$ to $-30$\,\%, displayed surface $-24$ to $-33$\,\% against each material's twin (one
seed each). Under weak gravity with a floor ($G=\rho g H/E=0.02$) every dragon run reaches its target (particle IoU
0.9870--0.9887 against the sample lifted onto the floor), and its centre of mass falls and comes to rest at the
floor-standing target's height ($-0.072$ to $-0.085$ world units against $-0.078$); the render term's lead under gravity is
$-71$\,\% on the exterior silhouette. On the bunny under the same load, two of three seeds and the twin stop after seven
windows: the landing raises the selection merit by 7--9\,\%, above the acceptance rule's 5\,\% bound, and three
rejections in a row end the run (the third seed lands and reaches its target, particle IoU 0.9863). At 40k a stronger
gravity, $G=0.05$, stops three of four targets early.
% TODO(results pending: single-body gravity through the multi-body scene path, with each body's objective read in a
% standing frame against a world-fixed target; the 40k stage has passed and the 300k runs are going. Replace the two
% sentences above with the 300k result.)

\subsection{Ablations}
\label{ch4:sec_ablation}
Table~\ref{ch4:tab_ablation} switches off one part at a time at 300k (bunny and dragon, three seeds each, against the
full method of the same seeds).

\begin{table}[t]
\centering
\footnotesize
\setlength{\tabcolsep}{4pt}
\caption{Ablations at 300k (bunny / dragon, three seeds; changes against the full method, all with the seed ranges
apart unless marked ``within''). Exterior residuals against an independent sample.}
\label{ch4:tab_ablation}
\begin{tabular}{p{3.6cm}p{9.8cm}}
\toprule
part switched off & what changes \\
\midrule
layer offset $u$ & exterior silhouette $\times3.0$ / $\times7.7$, shading $\times2.2$ / $\times4.0$, displayed surface
  $\times2.0$ / $\times3.4$, thin $+26\,\%$ / $\times3.0$; never at rest (end kinetic $\times54$ / $\times53$); early
  floating particles $-99$ / $-98\,\%$ \\
relief reference of the relaxation & dragon: exterior silhouette $+16\,\%$, shading $+20\,\%$; bunny thin $-23\,\%$ \\
volume-conserving body (carried volume, density match, volumetric assimilation) & early floating particles $\times5.0$ /
  $\times3.4$, end kinetic $+79$ / $+69\,\%$, dragon rotation $\times2.0$; in the kept dragon run, pieces of up to 42
  particles reach 20 pitches from the body at 0.4--3.6\,s; thin $-25$ / $-21\,\%$ \\
\quad carried volume alone & early floating particles $\times5.1$ / $\times3.5$ \\
\quad volumetric assimilation alone & not at rest: end kinetic $\times36$ / $\times23$ \\
\quad density match alone & within the full method's range \\
render on the exterior (read on the particle cloud instead) & exterior silhouette $+28$ / $+44\,\%$, shading $+21$ /
  $+19\,\%$, displayed surface $+27$ / $+20\,\%$; particle IoU and thin within \\
render weight averaged over windows & render share 0.42--0.80 in windows 3--11; end within the range but for small
  shading and thin changes \\
minimum spacing & surface roughness $\times2.1$, thin $\times2.9$ / $\times3.8$, exterior silhouette $+37$ / $+41\,\%$,
  chamfer $+5$ / $+6\,\%$, visible pieces $\times6.9$ / $\times7.6$ \\
relaxation in the driven half only (in every step instead) & late wobble returns: in-window normal excursion
  $\times2.1$ / $\times1.8$; thin $+14$ / $+30\,\%$, shading $+7$ / $+13\,\%$; runs end sooner \\
carried moments of $u$ (fresh every window instead) & window-to-window reversal returns (cosine $-0.47$ / $-0.31$ at
  7--15\,s against $-0.01$ / $+0.17$); visible pieces $\times2.0$ / $\times2.3$ \\
symmetric minimum-spacing pairs & bunny centre-of-mass drift returns ($1.0$--$1.3\times10^{-4}$ against
  $2.8$--$4.0\times10^{-5}$ pitches); shape within a few per cent \\
MPM cell following $N$ (instead of the shape's) & mixed: gaps open earlier and thin parts improve on some meshes, while
  bob's exterior silhouette is $+34\,\%$; bunny and bob end less at rest, armadillo and homer more \\
\bottomrule
\end{tabular}
\end{table}

Three readings stand out. First, $u$ is the main shape lever: without it every surface measure is two to eight times
worse and the body never comes to rest, while the early floating particles almost vanish, because $u$ is also their
source in the first seconds. Second, the exterior read is ahead of the cloud read on every measure of the surface,
while the cloud read is ahead on the particle measures it optimises; at 40k the cloud read's particle IoU is ahead on 17
of 19 meshes. Each read wins on the representation it scores, and the exterior is the one a viewer sees. Third, each of
the three late refinements earns its place on what it was made for (the wobble, the reversal and the centre of mass)
at a cost in run length on the dragon. We call these three parts (the driven-only
relaxation, the carried moments of $u$ and the symmetric spacing pairs) the late refinements. At 40k, the driven-only relaxation also improves the particle silhouette on 13 of
19 meshes, while the carried moments and the symmetric spacing are neutral there.

\subsection{Stress control against force control}
Classical physics-based shape control drives matter toward a target with external
forces~\citep{treuille2003keyframe,mcnamara2004adjoint,fattal2004targetdriven,chen2024eigenfluids,li2025flowmapadjoint}.
To test the choice of control, we replaced $\delta\mathbf{F}_c$ by an optimised external force per particle and driven
step, keeping everything else (objective, render terms, $u$, optimiser, windows) unchanged. A second arm projects the
force at every step onto zero net force and zero net torque about the centre of mass, the natural momentum-free
alternative. On six meshes and three seeds:
\begin{itemize}
  \item The plain force puts momentum in at every window: the centre of mass wanders 0.1--2.9 pitches and the body turns by
        up to $3.9^\circ$, $10^3$--$10^4$ and $10^2$--$10^3$ times the stress control's, with 77--158 times its angular
        momentum, and the objective has to pull it back.
  \item The projected force removes the net momentum exactly, but not the shape deficit. Both force arms end with
        3.1--5.7 times (plain) and 4.5--5.1 times (projected; bob 207 times) the stress control's geometry energy,
        17--71\,\% and 36--97\,\% more exterior silhouette error, 15--62\,\% more shading error, a 0.3--3.9\,\% worse chamfer
        distance and 1.5--25 times the surface fragments at the end, with the seed ranges apart in all but one case,
        at the same render share.
  \item The projected force morphs 1.6--2.4 times slower and on bob stalls with a lifted skin.
\end{itemize}
The force arms do not end less at rest, so the advantage of the stress control is not a spring-back argument: it is exact
momentum without a projection and a better end that a momentum-free force does not reach with the same objective.

\subsection{Comparison with other methods}
\label{ch4:sec_comparison}
We compare with nine methods of five kinds, each run from the same sphere toward the same six targets and read by the
same probes (Table~\ref{ch4:tab_methods}). The frames of mesh- and SDF-based methods are sampled by our volume sampler
at 300k in the target's frame; particle methods keep their particles. Thin coverage is read here on one set of thin
points per mesh, taken from the independent sample and the same for every method (it differs from the run's own thin
measure used above).
\begin{itemize}
  \item \textbf{Chapter~\ref{CH3_SHAPEMORPHING}} (our earlier F-control formulation~\citep{11088224}), reimplemented in our
        simulator and verified at the published scale ($N=11{,}160$: $-97.9$\,\% loss reduction against the published
        $-98.3$\,\%). It keeps the published loss, schedule and morph lengths, with the step size chosen by its own loss.
  \item \textbf{Neural implicit morphs.} The volume-preserving neural shape morph~\citep{buonomo2025vpnsm} and
        NISE~\citep{novello2023nise}, both from the former's code with the authors' defaults; the implicit surface
        deformation with explicit velocity fields~\citep{sang2025isd} (ISD) and its follow-up 4Deform~\citep{sang2025fourdeform},
        with the authors' code and shipped configurations. ISD and 4Deform expect point correspondences, which a
        sphere-to-target pair does not have, so their matching weight is zero; a second arm gives them radial
        correspondences (each sphere point paired with the ray hit from the centre) on the bunny and the dragon.
  \item \textbf{As-rigid-as-possible interpolation}~\citep{alexa2000arap} in gradient-domain
        form~\citep{sumner2004deftransfer,xu2005poisson}, with the sphere correspondence from conformalized mean curvature
        flow~\citep{kazhdan2012cmcf}; genus-zero targets only.
  \item \textbf{Optimal-transport displacement interpolation} on our particles~\citep{mccann1997convexity,bonneel2011displacement}:
        transport without a body.
  \item \textbf{Render-only mesh optimisation} from the sphere with our image model: Large
        Steps~\citep{nicolet2021largesteps} and FlexiCubes~\citep{shen2023flexicubes}, each with our 18 views at 188\,px
        (a) and with 25 views at 512\,px (b).
\end{itemize}

\begin{table*}[t]
\centering
\footnotesize
\setlength{\tabcolsep}{3pt}
\caption{Every method on the six 300k targets (one reader for all; per mesh the median over seeds, then the range over
the meshes the method ran on). Silhouette IoU and chamfer against an independent sample, also as ratios to the sampling
floor (a second sample of the mesh; for the thin points the floor is 2.9--8.1\,\%). ``Ours ahead'': the meshes on which
our seed range lies above the method's; no method's range lies above ours on any mesh. Volume: one voxel-occupancy
routine over every frame. $^a$Counted over the closed ends; the ends are open on two (neural morph) and three (NISE)
meshes. $^b$Plus 20--39\,min of signed distance fitting per shape. $^c$4Deform stays near the sphere, which matches only
the genus-zero targets. $^d$Genus-zero targets only.}
\label{ch4:tab_methods}
\resizebox{\textwidth}{!}{%
\begin{tabular}{lcccccccccc}
\toprule
method & meshes & silhouette IoU & $1-\mathrm{IoU}$ & chamfer & thin uncovered & ours ahead & end genus & volume over & frames with a piece & wall \\
 & & (indep.) & / floor & / floor & (\%) & (meshes) & = target & the frames & $\geq5$ pitches off & time \\
\midrule
ours & 6 & 0.9864--0.9899 & 1.1--1.3 & \textit{pending} % TODO(results pending: the one-sampler chamfer of our display surface)
 & 4.1--6.1 & -- & \textit{pending} % TODO(results pending: genus and watertightness of our display surface)
 & 0.96--1.09 & 0--8\,\% & 8.8--33.7\,min \\
Chapter~\ref{CH3_SHAPEMORPHING} & 6 & 0.9396--0.9754 & 2.2--7.4 & 1.03--2.20 & 30.2--48.4 & 6 of 6 & (particles) & 0.72--1.35 & 7--98\,\% & 2.3--6.8\,min \\
vol.-preserving neural morph & 6 & 0.9163--0.9688 & 3.8--9.2 & 1.05--1.66 & 29.8--45.1 & 6 of 6 & 0 of 4$^a$ & 0.73--1.47 & 0--71\,\% & 45--47\,min$^b$ \\
NISE & 6 & 0.9660--0.9830 & 1.9--3.0 & 1.01--1.02 & 20.1--35.2 & 6 of 6 & 3 of 3$^a$ & 0.67--1.01 & 0--24\,\% & 13\,min$^b$ \\
ISD & 6 & 0.9822--0.9901 & 1.2--1.6 & 1.00--1.01 & 12.2--18.1 & 3 of 6 & 5 of 6 & 0.00--5.09 & 5--38\,\% & 27--39\,min \\
ISD, radial corr. & 2 & 0.9816--0.9862 & 1.3--1.7 & 1.01 & 14.3--19.0 & 1 of 2 & 1 of 2 & 0.36--3.99 & -- & 27--34\,min \\
4Deform & 6 & 0.4955--0.7628 & 26--57 & 2.8--10.3 & 67.5--81.5 & 6 of 6 & 3 of 6$^c$ & 0.24--2.10 & 0--10\,\% & 68--138\,min \\
4Deform, radial corr. & 2 & 0.9329--0.9784 & 2.0--6.0 & 1.04--3.94 & 22.7--32.2 & 2 of 2 & 0 of 2 & 0.86--1.48 & -- & 69--75\,min \\
ARAP$^d$ & 3 & 0.9849--0.9884 & 1.2--1.3 & 1.01 & 9.7--16.8 & 1 of 3 & 3 of 3 & 0.86--1.02 & 0\,\% & 0.2--12\,min \\
OT displacement & 2 & 0.9306--0.9440 & 5.0--6.5 & 1.04--1.11 & 24.0--24.7 & 2 of 2 & (particles) & 0.96--1.28 & 76--81\,\% & -- \\
Large Steps (a) & 6 & 0.9777--0.9881 & 1.2--2.0 & 1.01--1.46 & 9.7--17.7 & 3 of 6 & 3 of 6 & 0.85--1.74 & 0\,\% & 18--33\,s \\
Large Steps (b) & 6 & 0.9807--0.9897 & 1.2--2.1 & 1.01--1.17 & 9.6--16.3 & 3 of 6 & 3 of 6 & 0.84--1.62 & 0\,\% & 27--53\,s \\
FlexiCubes (a) & 6 & 0.8307--0.9315 & 7.5--20.7 & 1.21--6.64 & 53.5--70.8 & 6 of 6 & 0 of 6 & 0.76--1.10 & 67--86\,\% & 10--21\,s \\
FlexiCubes (b) & 6 & 0.9055--0.9803 & 2.2--11.6 & 1.02--3.61 & 31.2--60.8 & 6 of 6 & 0 of 6 & 0.73--1.20 & 57--90\,\% & 17--28\,s \\
\bottomrule
\end{tabular}}
\end{table*}

\paragraph{The end.} No method's silhouette IoU range lies above ours on any mesh. Ours lies above every other method,
with the seed ranges apart, on the dragon, the teapot and the armadillo, and is level with the best three on the other
three meshes: on the bunny with ISD's radial arm (0.9862), Large Steps (0.9868--0.9875) and ARAP (0.9874) against our
0.9860--0.9891; on bob with ISD (0.9894--0.9911) and Large Steps (b) (0.9893--0.9899) against 0.9883--0.9908; on homer
with ISD (0.9825--0.9878), Large Steps (b) (0.9880--0.9888) and ARAP (0.9884) against 0.9869--0.9890. On the shared thin
points, ours leaves 4.1--6.1\,\% uncovered, as many as a second sample of the mesh (2.9--8.1\,\%); every other method
leaves at least 9.2\,\% (Large Steps 9.2--16.4, ARAP 9.7--16.8, ISD 10.5--22.1, NISE 20--35, the volume-preserving morph
26--62, Chapter~\ref{CH3_SHAPEMORPHING}'s formulation 29--49, FlexiCubes 29--73, 4Deform 30--98\,\%). The margins are
widest on the surface measures of this chapter, where the render terms act. Where they were read, the bunny's exterior
silhouette residual is $7.0$--$7.7\times10^{-4}$ against ARAP's $10.7$, ISD's $20.5$--$20.8$,
Chapter~\ref{CH3_SHAPEMORPHING}'s $90$--$99$ and the volume-preserving morph's $106$ (all $\times10^{-4}$), and the dragon's
$5.0$--$5.8$ against ISD's $31$--$45$ and Chapter~\ref{CH3_SHAPEMORPHING}'s $68$--$71$; ARAP is level with ours on the
displayed surface of the armadillo and homer and behind on their exterior silhouette. The formulation of
Chapter~\ref{CH3_SHAPEMORPHING} also does not come to rest (end kinetic measure 0.0028--0.95, two to seven orders of
magnitude above ours), since its objective scores only the mass on the grid.

\paragraph{The way there.} What separates the methods most is what happens between the sphere and the target.
\begin{itemize}
  \item \textbf{ISD} fits the end but does not keep a body on the way. In every ISD run the volume collapses in mid-run, to
        0.1--65\,\% of the target's (0.1--0.3\,\% on bob, the dragon and homer), and on bob and the armadillo it swells to
        five times the target's before ending at 0.98--1.03. Between 6 and 21 of its 21 frames have more than one
        component, and its transport progress jumps past 0.25, 0.5 and 0.75 within one frame on five meshes.
  \item \textbf{Large Steps} keeps one genus-zero surface in every frame. It never opens the dragon's fifteen handles, the
        teapot's handle or bob's ring and imitates them with self-intersecting surfaces (31 of 36 ends self-intersect),
        and its volume swings to 1.6--1.7 times the target's on the dragon.
  \item \textbf{FlexiCubes} changes topology without control: spurious genus 4--58 at the end (up to 68 over the frames)
        on targets of genus 0--15, with up to 149 components.
  \item \textbf{4Deform} without correspondences does not leave the sphere (silhouette IoU 0.50--0.76, its transport
        progress at the end at most 0.28); given radial correspondences, it reaches 0.978 on the bunny and 0.933 on the
        dragon.
  \item \textbf{The volume-preserving neural morph and NISE} end with open surfaces on two and three meshes, and five runs
        of the former hold a static field whose zero set contains both shapes at every time (the teapot's three seeds
        among them). NISE is 3.4 times faster than the volume-preserving morph and ends above it on all six meshes.
  \item \textbf{Transport alone and Chapter~\ref{CH3_SHAPEMORPHING}'s formulation} keep the mass but tear thin parts into
        clouds: pieces farther than five pitches from the body appear in 76--81\,\% and 7--98\,\% of their frames, and
        at transport progress $p=0.5$ (the share of the transport distance to the independent sample already covered)
        1543--4042 and 788--45009 particles lie outside the largest connected piece.
\end{itemize}
Ours stays one body with exact mass. Its physical volume stays within $-0.8$ to $+3.1$\,\% of the target's, the common
occupancy routine reads 0.96--1.09 over every frame (the high end is the dragon's early transient), pieces farther than
five pitches occur in 0--8\,\% of the frames with at most 14 particles in them, and it reaches every matched-progress
point on all six meshes. At $p=0.5$ it has 0--6 particles outside the largest connected piece on four meshes, 93 on the
bunny (Large Steps 60--85, ARAP 59) and 2622 on the teapot (Large Steps 33--69). In a count over the first seconds on the
bunny and the dragon, about nine in ten such particles were bumps about half a pitch high on the surface rather than
separate pieces.

\paragraph{Cost.} On one RTX~6000 Ada, Large Steps (18--53\,s), FlexiCubes (10--28\,s), ARAP (12\,s--12\,min),
Chapter~\ref{CH3_SHAPEMORPHING}'s formulation (2.3--6.8\,min) and NISE (13\,min) are faster than ours (8.8--33.7\,min);
ISD (27--39\,min), the volume-preserving morph (45--47\,min plus its signed distance fits) and 4Deform (68--138\,min) are
slower. In short, no other method is at once a body with exact mass, reaches every target including those of higher
genus, and keeps thin parts at the sampling floor; the render terms carry the widest margins on the surfaces they read
(1.4 to 600 times on the exterior silhouette) and the narrowest on the particle silhouette IoU, which they do not move.

\subsection{Runtime}
\label{ch4:sec_runtime}
\begin{table}[t]
\centering
\footnotesize
\caption{Wall-clock time of a 300k morph on one RTX 6000 Ada (full method; ranges over seeds) and its committed windows
(median).}
\label{ch4:tab_runtime}
\begin{tabular}{lcc}
\toprule
mesh & minutes & windows \\
\midrule
bunny & 8.8--11.3 & 87 \\
dragon & 23.3--28.7 & 165 \\
teapot & 12.6--13.7 & 123 \\
bob & 15.5--16.1 & 130 \\
armadillo & 12.3--14.8 & 90 \\
homer & 23.4--33.8 & 197 \\
\bottomrule
\end{tabular}
\end{table}
Table~\ref{ch4:tab_runtime} lists the wall-clock times; the card's peak memory is 19.3--23.9\,GB. An optimiser attempt
takes 6.5\,s on the bunny and 8.8\,s on the dragon. Early in a dragon run the Sinkhorn solves take 32\,\% of an attempt,
the surface proximity 17\,\% and the three adjoint sweeps 18\,\%; late, the sweeps take 28\,\% and the Sinkhorn solves
21\,\%. The two exact speed-ups of Section~\ref{sec:implementation} save 4.4--6.1\,\% a window. The run length is set by
the stop rule: within one seed standard deviation, the end measures settle 0.6--1.2\,min before the end on the bunny and
1.2--4.2\,min on the dragon, and no stop rule read from the run's own noise saves more than about a minute, because the
objective really keeps improving. Stopping at a fixed 5\,s would take 4.0\,min on the bunny and 7.2\,min on the dragon,
but at 5\,s the exterior silhouette residual is still 1.19 and 2.44 times its end value and neither body is at rest.

\subsection{Two thrown bodies}
\label{ch4:sec_throw}
Two spheres of 150k particles each are thrown toward each other, collide in the air while half-formed, part, and each
becomes a bunny that lands and tumbles (Figure~\ref{ch4:fig_throw}). The scene uses the multi-body definitions of
Section~\ref{ch4:sec_scene} and the stiff material ($E\times7$), with the gravity and throw of a scene designed for the
default material at $G=0.02$, re-expressed so that the motion in world units is unchanged. The plastic assimilation is
off in the scene, so the controls hold each body's shape against the impacts. With the default material the bunnies
flattened at the landing and their ear tips were drawn out two- to sixfold on average; the stiffer material meets the
floor at a third of the default material's impact speed relative to the wave speed (0.12 against 0.31) and holds its
shape.

In the delivered run the bodies first touch at 3.8\,s, when the morph is 43\,\% complete (silhouette progress in each
body's frame), part in the air at 5.7\,s with a restitution of 0.95, land at 8.95 and 10.95\,s, spring off the floor and
tumble by $68^\circ$ and $127^\circ$. At the end each body's silhouette IoU against its target in its own frame is 0.985
and 0.984 (end shape error 0.015), 4.9 and 3.3\,\% of the thin points are uncovered, the ear tips stretch by 23 and 15\,\%
on average through the landing, and no piece of ten or more particles leaves a body by more than three spacings at any
time. Momentum is kept within $2.7\times10^{-4}$ of the bodies' momentum in every phase. The bodies are still moving at the
end: they are tumbling. The render-off twin passes the same checks with one four-particle piece; a second seed lands one
body on its head, which deforms the head band by about nine spacings for a second before it recovers, and misses the
angular momentum check by $1.05\times10^{-3}$ against $10^{-3}$. The run takes 50.5\,min. The render term's lead in this
scene (end silhouette 0.985 / 0.984 against 0.983 / 0.984) is about the size of the seed's spread.
% TODO(results pending: the 300k showcase run that has just started; add its result and frames here once read.)

\begin{figure*}[t]
    \centering
    \includegraphics[width=\textwidth]{figs/ch4/fig_throw_scene.jpg}
    \caption{Two thrown bodies (2 $\times$ 150k particles, stiff material). Left to right, top to bottom: launch,
    $t=2$\,s, first contact (3.85\,s), parted (5.7\,s), 8\,s, after the first landing (9.2\,s), after the second landing
    (11.2\,s), and the end (16\,s). Each body has its own grid field and its own objective in its own frame, so the
    bodies part cleanly, fly, land and tumble while each becomes a bunny.}
    \label{ch4:fig_throw}
\end{figure*}

% ============================================================================
\section{Measured Costs and What Was Tried}
\label{ch4:sec_limitations}
% ============================================================================

Four costs of the method are measured, and each has been attacked; we report what the attempts did.

\begin{itemize}
  \item \textbf{Small transient floaters.} The early arrival (windows 3--10) leaves particles apart from the body,
        mostly relief standing half a pitch outside the displayed surface; a viewer sees a few separate pieces, 2--38
        particles in all over windows 3--10 (side view, medians over seeds). None persists and none is serious by our
        definition. $u$ is their main source: without it they almost vanish, and the surface measures are two to eight
        times worse. Keeping each layer step within the sampling berth of its nearest neighbour removed three quarters of
        the bunny's visible pieces, but not the dragon's early ones, and it moved the centre of mass at every step (the
        rigid part of the steps not taken), so it was not adopted.
  \item \textbf{Narrow gaps.} The cubic kernel averages the grid velocity over about two cells (twelve particle spacings at
        300k). A gap of the target narrower than about two cells shares grid nodes with both rims, cannot open until late,
        and tears into beads there (the dragon's gap between hind leg and body at about 3--6\,s). An MPM cell that follows
        $N$ opens these gaps earlier but prints surface features ahead of the bulk, and neither cell wins on every mesh
        (Table~\ref{ch4:tab_ablation}). Putting the two sides of each target gap on their own grid fields, with free
        separation and no tension across the gap, removed the crevice's beads at 3--6\,s in every dragon run, but new
        beads of one to three particles appeared where the gaps then opened, earlier, and it was not adopted.
  \item \textbf{Runtime.} A 300k morph takes 9--34 minutes: slower than Large Steps, FlexiCubes, ARAP,
        Chapter~\ref{CH3_SHAPEMORPHING}'s formulation and NISE, faster than ISD, the volume-preserving neural morph and
        4Deform (Section~\ref{ch4:sec_comparison}). The exact speed-ups (cell-ordered transfers, skipped
        idle layer kernels, reused neighbour lists) are in the method; the next exact candidates are each below 2\,\% a
        window. The rest is the run length, which the objective's own progress sets (Section~\ref{ch4:sec_runtime}).
  \item \textbf{Angular momentum and the drag.} The frozen drag damps the particle momentum but not APIC's affine part,
        so it is not angular-momentum consistent, and its effect grows with the motion (bob's rotation and the angular
        momentum differences above). A drag that damps the whole APIC momentum cut the largest angular momentum about
        fourfold on the 40k bunny, kept shape and rest within the seed spread, and let the body turn about seven times
        more from shape change, which the frozen drag had been suppressing; it is not part of the frozen method.
\end{itemize}
% TODO(results pending: the late tail. A target that moves along the transport geodesic over the body's elastic
% period spreads the deformation in time (40k: 81--91 % of the movement after 1 s, against 14--60 %). With u gated on the
% arrival at the final target, the end came mostly back but stayed a few per cent short on two or three meshes, because
% the interpolated in-between targets ask for 3--14 % more volume; the next form rescales them to the body's density.
% Write the outcome here, as a result, once read.)
% TODO(results pending: the morph toward a target given only by images, the render terms then the only shape signal;
% add it as a subsection of the results once read.)

% ============================================================================
\section{Summary}
\label{ch4:sec_summary}
% ============================================================================

This chapter turned the render-guided idea of our earlier formulation into a method whose every part is measured. The
body is a volume-conserving MLS-MPM solid driven by a stress increment and a rigid-free surface offset; it is scored in
settled windows by a transport energy, two settling terms, two local cleanup terms and the silhouette and shading of its
own exterior surface; and the render term holds a third of each window's first gradient. Against its $\lambda=0$ twin the render term
lowers the error of the surface a viewer sees by a fifth to a half on six meshes, with an outline closer to the exact
mesh than a perfect particle sample, and a causal switch shows that it maintains that surface rather than storing it. The morph keeps
exact mass and, up to the ambient drag, momentum, ends at rest, opens handles a sphere does not have, and takes part in
contact and gravity with other bodies. Its measured costs, small transient floaters, late-opening narrow gaps and a run
time of minutes, are reported with what was tried against each.
